% \section{Abstract}
% Prior work on document generation has generally tackled the creation of each separate format to be a different task, leading to fragmented learning processes, redundancy in models and methods, and disjointed evaluation. We consider each of these documents as \textit{templatic views} of the same underlying knowledge/content, and we aim to unify the generation and evaluation of these templatic views. We begin by showing that current LLMs are capable of generating various document formats with little to no supervision. Further, a simple augmentation involving a structured intermediate representation can improve performance, especially for smaller models. We then introduce a novel unified evaluation framework that can be adapted to measuring the quality of document generators for heterogeneous downstream applications. This evaluation is adaptable to a range of user defined criteria and application scenarios, obviating the need for task specific evaluation metrics. Finally, we conduct a human evaluation, which shows that people prefer 82\% of the documents generated with our method, while correlating more highly with our unified evaluation framework than prior metrics in the literature.


\begin{figure*}[!h]
    \centering
    \includegraphics[width=\textwidth]{research-chapters/3-chapter/images/sure_method_horizonal.pdf}
    \caption[Visualization of our method to unify the generation and evaluation of templatic views of documents.]{Visualization of our method to unify the generation and evaluation of templatic views of documents. Given an input document, we prompt the LLM to generate an intermediate representation. We can use the representation to prompt the model to generate a templatic view of the input document. We then evaluate the generations using our unified evaluation framework. The LLM represented in the figure is the same model.}
    \label{fig:method}
    \vspace{-2mm}
\end{figure*}

\section*{Non-technical abstract}
We often need to repackage the same information into different formats for different audiences. A product manager might turn a single product proposal into a slide deck, a blog post, and a requirements document. A scientist might transform one research project into a conference talk, a poster, a blog post, and a series of social media updates. Each of these is essentially a summary, just shaped to fit a specific format and audience.
In this chapter, we show that AI systems can produce better versions of these documents by working in two steps instead of one. First, the AI extracts the key information from the source document and organizes it into a structured outline. Then, it uses that outline as a blueprint to write the final document in whatever format is needed. This two-step approach noticeably improves quality, and the benefits are especially pronounced for smaller, cheaper AI models, making them more competitive with their more expensive counterparts.
We also developed a new way to automatically score these kinds of summaries that takes their intended structure into account. Compared to existing scoring methods, ours lines up much more closely with how human readers actually judge quality.


\section{Introduction}
Sharing information is vital for communication and discourse across domains, as it allows for knowledge to be disseminated to a wider audience. This is often done by users through documents in multiple formats that nevertheless share some underlying knowledge. A product manager may need to create a requirements spec, a product pitch deck, and an announcement newsletter for the same project. Likewise, a person on the job market may create a resume, a cover letter, and a personal website. We consider these documents to be \textit{templatic views} of the same underlying knowledge.

This is equally true for the scientific domain, in which researchers create documents in multiple formats to effectively communicate and showcase their work -- such as through academic papers, conference talks, social media posts, poster presentations, and non-technical blog posts. Sharing knowledge in multiple formats broadens the audience and can help bridge the information gap between domain experts, researchers in adjacent fields, and even the general public, leading to greater understanding, collaborations and accelerated progress~\cite{Bornmann2014GrowthRO}.


Past work on document generation has focused on developing generation and evaluation methods specific to a single document type~\cite{Fu2021DOC2PPTAP,Qiang2016LearningTG,chandrasekaran2020overview}. 
Narrow, custom methods tailored to individual document types are time consuming to engineer and manage over the long term. For example, in an enterprise setting, it’s common to have dozens of occupation- and task-specific documents, each with their own template.
Additionally, specific trained methods require data that may be expensive to acquire, or even be unavailable entirely. Meanwhile, LLMs have recently shown great success in long document generation~\cite{radford2019language,Brown2020LanguageMA}, indicating that this fragmentation of methods may no longer be necessary. Thus, our goal is to unify methods for both generating and evaluating templatic views of documents, allowing system designers and engineers to manage and adapt to a range of document types and domains easily and efficiently. 

We begin by showing that LLMs are capable of diverse, structured document generation, requiring very little instructional guidance to do so effectively. Additionally, a few minor augmentations to the prompt -- such as a structured, intermediate representation, and simple stylistic descriptions -- can further improve downstream performance, especially for smaller, less resource intensive models. These findings have important implications on the deployment and scaling of unified, real-world AI-assisted document authoring systems.

We then introduce Template Adaptable Evaluation (TAE), departing from prior work's task specific evaluation methods~\cite{zhang2020bertscore,Qiang2016LearningTG,Wang2015iPosterIP}. TAE is a unified precision-recall style framework for automatic evaluation that is highly customizable, allowing users to easily integrate existing text-based metrics from the literature into its formulation and tailor it to their specific use case. Additionally, this framework allows developers to compare performance across document types, without needing to develop an evaluation metric for each individual template.


We evaluate our unified approach for templatic view generation and evaluation on 3 types of documents: slides, posters, and blog posts~\cite{Fu2021DOC2PPTAP,Qiang2016LearningTG,chandrasekaran2020overview}. %We use the DOC2PPT~\cite{Fu2021DOC2PPTAP}, Paper-Poster~\cite{Qiang2016LearningTG}, and LongSumm~\cite{chandrasekaran2020overview} datasets, respectively, where the inputs are scientific articles and the outputs are the documents of the corresponding types.
Our experiments demonstrate that using a structured intermediate representation leads to improvements in performance across tasks, with greater gains for smaller language models.
%demonstrating that LLMs are capable of generating templatic views with structure-augmented prompts, despite having no supervision.
In our human evaluation to validate both our unified document generation method and evaluation metric, we show that annotators prefer the output yielded by the structure-aware generation process 82\% of the time and that our evaluation metric correlates more highly with human preference than other popular metrics. A visualization of our methods can be found in Figure~\ref{fig:method}. 
